% Guide utilisateur VetNutri MP — version desktop
% Compilation : latexmk -pdf guide_utilisateur.tex
% Les captures (dossier captures/) sont produites par le pilote
% composeApp/src/desktopTest/.../Guide/GuideScreenshotDriver.kt (voir README.md).
% Les repères numérotés sont placés en pixels de capture (1400 x 900) : \repere{x}{y}{n}.

\documentclass[11pt,a4paper]{article}

\usepackage[utf8]{inputenc}
\usepackage[T1]{fontenc}
\usepackage{lmodern}
\usepackage[french]{babel}
\usepackage[margin=2.2cm]{geometry}
\usepackage{graphicx}
\usepackage{float}
\usepackage[font=small,labelfont=bf]{caption}
\usepackage{xcolor}
\usepackage{enumitem}
\usepackage{fancyhdr}
\usepackage{titlesec}
\usepackage{amsmath}
\usepackage{array}
\usepackage{booktabs}
\usepackage{tabularx}
\usepackage{tikz}
\usetikzlibrary{positioning,arrows.meta}
\usepackage[most]{tcolorbox}
\usepackage[hidelinks]{hyperref}

\definecolor{vetprimaire}{HTML}{00897B}
\definecolor{vetfond}{HTML}{E0F2F1}
\definecolor{vetalerte}{HTML}{B00020}
\definecolor{vetrepere}{HTML}{E65100}
\definecolor{vetinfo}{HTML}{1565C0}

\titleformat{\section}{\Large\bfseries\color{vetprimaire}}{\thesection}{0.6em}{}
\titleformat{\subsection}{\large\bfseries\color{vetprimaire}}{\thesubsection}{0.6em}{}
\titleformat{\subsubsection}{\normalsize\bfseries}{}{0em}{}

\pagestyle{fancy}
\fancyhf{}
\fancyhead[L]{\small VetNutri MP — Guide utilisateur}
\fancyhead[R]{\small Version desktop 3.3.18}
\fancyfoot[C]{\small\thepage}
\renewcommand{\headrulewidth}{0.4pt}

\setlist{itemsep=2pt,topsep=4pt}
\setlength{\parindent}{0pt}
\setlength{\parskip}{4pt}

% ------------------------------------------------------------------ Encadrés
\newtcolorbox{aretenir}{colback=vetfond,colframe=vetprimaire,boxrule=0.6pt,arc=2pt,
  left=6pt,right=6pt,top=4pt,bottom=4pt,fontupper=\small,title=À retenir,fonttitle=\small\bfseries}
\newtcolorbox{pourquoi}{colback=blue!3,colframe=vetinfo,boxrule=0.6pt,arc=2pt,
  left=6pt,right=6pt,top=4pt,bottom=4pt,fontupper=\small,title=Pourquoi ?,fonttitle=\small\bfseries}
\newtcolorbox{attention}{colback=red!4,colframe=vetalerte,boxrule=0.6pt,arc=2pt,
  left=6pt,right=6pt,top=4pt,bottom=4pt,fontupper=\small,title=Attention,fonttitle=\small\bfseries}
\newtcolorbox{verifiez}{colback=green!4,colframe=green!45!black,boxrule=0.6pt,arc=2pt,
  left=6pt,right=6pt,top=4pt,bottom=4pt,fontupper=\small,title=Ce que vous devez voir,fonttitle=\small\bfseries}

% Étape du cas pratique : \etape{titre}{objectif}
\newcommand{\etape}[2]{%
  \subsection{#1}
  \begin{tcolorbox}[colback=white,colframe=vetprimaire!60,boxrule=0.4pt,arc=2pt,left=6pt,
    right=6pt,top=3pt,bottom=3pt,fontupper=\small]
  \textbf{Objectif :} #2
  \end{tcolorbox}}

% Repère numéroté sur une capture (coordonnées en pixels de l'image 1400 x 900)
\tikzset{repere/.style={circle,fill=vetrepere,draw=white,line width=0.8pt,text=white,
  font=\sffamily\bfseries\footnotesize,inner sep=0pt,minimum size=15pt}}
\newcommand{\repere}[3]{\node[repere] at ({#1/1400},{1-#2/900}) {#3};}
% Rappel d'un repère dans le texte
\newcommand{\rep}[1]{\tikz[baseline=(n.base)]\node[repere,minimum size=12pt,font=\sffamily\bfseries\scriptsize] (n) {#1};}

% Capture simple : \capture{fichier}{légende}{label}
\newcommand{\capture}[3]{%
  \begin{figure}[H]
    \centering
    \fbox{\includegraphics[width=0.92\linewidth]{captures/#1}}
    \caption{#2}
    \label{fig:#3}
  \end{figure}}

% Capture annotée : \captureannotee{fichier}{légende}{label}{repères}
\newcommand{\captureannotee}[4]{%
  \begin{figure}[H]
    \centering
    \fbox{\begin{tikzpicture}
      \node[anchor=south west,inner sep=0] (img)
        {\includegraphics[width=0.92\linewidth]{captures/#1}};
      \begin{scope}[x={(img.south east)},y={(img.north west)}]
        #4
      \end{scope}
    \end{tikzpicture}}
    \caption{#2}
    \label{fig:#3}
  \end{figure}}

% Libellé d'un bouton ou d'un élément de l'interface
\newcommand{\ui}[1]{\textsf{\textbf{#1}}}

\begin{document}

% ================================================================== Titre
\begin{titlepage}
  \centering
  \vspace*{2cm}
  \includegraphics[width=4cm]{logo.png}\par\vspace{1cm}
  {\Huge\bfseries\color{vetprimaire} VetNutri MP\par}
  \vspace{0.4cm}
  {\Large Gestionnaire de nutrition vétérinaire\par}
  \vspace{1.5cm}
  {\LARGE Guide utilisateur — version desktop\par}
  \vspace{0.3cm}
  {\large De la prise en main à la première ordonnance\par}
  \vspace{0.4cm}
  {\large Version de l'application : 3.3.18\par}
  \vspace{1.2cm}
  {\large Sébastien Lefebvre\par}
  \vfill
  {\small Les captures d'écran ont été réalisées sur une base de démonstration.
  Les animaux, propriétaires et valeurs saisies sont fictifs.\par}
  \vspace{0.5cm}
  {\large Octobre 2026\par}
\end{titlepage}

\tableofcontents
\newpage

% ================================================================== 1
\section{Avant de commencer}

\subsection{Ce que vous saurez faire}

À la fin de ce guide, vous saurez :
\begin{itemize}
  \item créer le dossier d'un animal et enregistrer une consultation ;
  \item choisir le référentiel nutritionnel adapté à l'animal ;
  \item composer une ration, lire son analyse et l'ajuster ;
  \item préparer une ordonnance et un compte rendu pour le propriétaire.
\end{itemize}
Le cœur du guide est un cas pratique suivi du début à la fin (section~\ref{sec:cas}) : Filou,
un Labrador en surpoids. Reproduisez-le une première fois sur votre poste : c'est le moyen le
plus rapide de prendre l'application en main.

\subsection{La démarche en un coup d'œil}

Toute formulation de ration suit les mêmes six étapes. Chacune s'appuie sur la précédente : la
ration est rattachée à une consultation, et l'analyse compare la ration aux besoins calculés
pour cette consultation.

\begin{figure}[H]
\centering
\begin{tikzpicture}[
  etape/.style={rectangle,rounded corners=3pt,draw=vetprimaire,fill=vetfond,
    text width=2.25cm,align=center,minimum height=1.35cm,font=\footnotesize},
  fleche/.style={-{Stealth[length=2.5mm]},thick,vetprimaire}]
  \node[etape] (a) {\textbf{1. Dossier}\\animal ou troupeau};
  \node[etape,right=0.2cm of a] (b) {\textbf{2. Consultation}\\poids, NEC};
  \node[etape,right=0.2cm of b] (c) {\textbf{3. Référentiel}\\besoins};
  \node[etape,right=0.2cm of c] (d) {\textbf{4. Ration}\\aliments, quantités};
  \node[etape,right=0.2cm of d] (e) {\textbf{5. Analyse}\\et ajustement};
  \node[etape,right=0.2cm of e] (f) {\textbf{6. Ordonnance}\\compte rendu};
  \draw[fleche] (a) -- (b);
  \draw[fleche] (b) -- (c);
  \draw[fleche] (c) -- (d);
  \draw[fleche] (d) -- (e);
  \draw[fleche] (e) -- (f);
  \draw[fleche,dashed] (e.south) to[bend left=35] node[below,font=\scriptsize,text=black]
    {on recommence jusqu'à une ration satisfaisante} (d.south);
\end{tikzpicture}
\caption{La démarche de formulation dans VetNutri MP.}
\label{fig:demarche}
\end{figure}

\subsection{Comment lire ce guide}

\begin{itemize}
  \item Les libellés de l'interface sont écrits \ui{comme ceci}.
  \item Sur les captures, les repères \rep{1} \rep{2} \rep{3} renvoient aux explications qui
        suivent la figure.
  \item Quatre types d'encadrés vous aident à comprendre et à vérifier :
\end{itemize}

\begin{aretenir}
L'essentiel à mémoriser.
\end{aretenir}
\begin{pourquoi}
Le raisonnement nutritionnel derrière ce que fait l'application.
\end{pourquoi}
\begin{verifiez}
Le résultat attendu à l'écran. Si vous voyez autre chose, relisez l'étape.
\end{verifiez}
\begin{attention}
Une action irréversible ou une erreur fréquente.
\end{attention}

VetNutri MP est une aide à la décision. Il ne remplace ni l'examen clinique ni le jugement du
vétérinaire, qui reste responsable des données saisies et des rations prescrites. Les données
sont enregistrées sur votre ordinateur et ne sont transmises que si vous les exportez ou les
partagez.

% ================================================================== 2
\section{Installation et premier lancement}
\label{sec:lancement}

Au tout premier démarrage, la base de l'application est vide. Trois actions, à ne faire
qu'une fois, la rendent utilisable.

\subsection{Installer les données}

\begin{enumerate}
  \item Lancez VetNutri MP. Une fenêtre \ui{Mise à jour des données} s'ouvre.
  \item Cliquez sur \ui{Installer la mise à jour}. L'installation dure environ une minute.
\end{enumerate}

\capture{01_mise_a_jour_donnees.png}{Proposition d'installation des données au premier lancement.}{maj}

\begin{pourquoi}
Les données intégrées contiennent la base d'aliments, les référentiels de besoins et les
conseils. Sans elles, vous ne pourriez ni choisir de référentiel ni composer de ration. La même
fenêtre réapparaît quand une nouvelle version apporte des données plus récentes ;
\ui{Plus tard} reporte l'installation.
\end{pourquoi}

\subsection{Accepter les conditions d'utilisation}

\begin{enumerate}[resume]
  \item En bas de l'écran d'accueil, cliquez sur \ui{Lire et accepter les conditions}.
  \item Lisez les conditions, puis cliquez sur \ui{J'accepte les conditions}.
\end{enumerate}

\capture{03_cgu.png}{Conditions générales d'utilisation.}{cgu}

\subsection{Entrer dans l'application}

\begin{enumerate}[resume]
  \item Cliquez sur \ui{Continuer}.
\end{enumerate}

\captureannotee{04_accueil_continuer.png}{Écran d'accueil prêt à l'emploi.}{accueil}{%
  \repere{55}{484}{1}
  \repere{560}{756}{2}
  \repere{560}{816}{3}
  \repere{560}{872}{4}}

\rep{1} L'état de la base : nombre d'aliments, de référentiels et de conseils.
\rep{2} \ui{Continuer} ouvre l'application.
\rep{3} \ui{Restaurer une sauvegarde} recharge une sauvegarde (section~\ref{sec:sauvegardes}).
\rep{4} \ui{Mode examen} sert aux sessions d'examen encadrées.

\begin{verifiez}
L'encadré indique \ui{Base de données complète}, avec plus de 11\,000 aliments. Aux lancements
suivants, il suffira de cliquer sur \ui{Continuer}.
\end{verifiez}

% ================================================================== 3
\section{Découvrir l'écran principal}
\label{sec:liste}

L'écran principal est la liste des animaux. Tout part de là.

\captureannotee{05_liste_animaux.png}{La liste des animaux et ses commandes.}{liste}{%
  \repere{184}{84}{1}
  \repere{528}{84}{2}
  \repere{872}{84}{3}
  \repere{1215}{84}{4}
  \repere{40}{150}{5}
  \repere{1000}{150}{6}
  \repere{1358}{145}{7}
  \repere{1296}{258}{8}
  \repere{1356}{815}{9}
  \repere{1325}{32}{10}}

\begin{tabularx}{\linewidth}{@{}c X@{}}
\rep{1} & \ui{Liste des aliments} : consulter et modifier la base d'aliments (section~\ref{sec:aliments}). \\
\rep{2} & \ui{Données de calcul} : équations, bibliographie, référentiels, conseils (section~\ref{sec:donnees}). \\
\rep{3} & \ui{Exploration multiration} : calculer des rations en série (section~\ref{sec:multiration}). Ce bouton n'apparaît que si la fonction est activée. \\
\rep{4} & \ui{Import Rapide} : importer des dossiers d'animaux. \\
\rep{5} & Recherche par identifiant, nom, propriétaire ou race. \\
\rep{6} & Filtre par espèce. \\
\rep{7} & Tri et filtres avancés. \\
\rep{8} & La flèche ouvre la fiche de l'animal ; la corbeille rouge le supprime. \\
\rep{9} & \ui{+} crée un nouveau dossier. \\
\rep{10} & Paramètres de l'application (section~\ref{sec:parametres}). \\
\end{tabularx}

% ================================================================== 4
\section{Cas pratique : la ration de Filou, pas à pas}
\label{sec:cas}

\begin{tcolorbox}[colback=vetfond,colframe=vetprimaire,boxrule=0.6pt,arc=2pt,
  title=Le cas,fonttitle=\bfseries]
Filou est un Labrador retriever mâle entier de 7 ans. Il pèse 32~kg et sa note d'état
corporel (NEC) est de 7/9 : il est en surpoids. L'objectif est une perte de poids progressive
avec un aliment hypocalorique du commerce.

Nous allons créer son dossier, enregistrer la consultation, composer la ration, la corriger
grâce à l'analyse, puis préparer l'ordonnance.
\end{tcolorbox}

% ------------------------------------------------------------------
\etape{Étape 1 — Créer le dossier}{enregistrer l'identité de Filou.}

\begin{enumerate}
  \item Dans la liste des animaux, cliquez sur \ui{+} en bas à droite.
  \item Laissez \ui{Dossier : Individu} coché.
  \item Saisissez l'\ui{ID} (par exemple le numéro de dossier de la clinique), le \ui{Nom} et le
        \ui{Propriétaire}.
  \item Choisissez l'\ui{Espèce} (\emph{Chien}), puis saisissez la \ui{Race}, la
        \ui{Date de naissance} et le \ui{Sexe}.
  \item Validez avec le bouton coché en bas à droite.
\end{enumerate}

\captureannotee{06_creation_individu.png}{Le formulaire de création d'un individu.}{creation-individu}{%
  \repere{545}{152}{1}
  \repere{318}{465}{2}
  \repere{850}{697}{3}
  \repere{318}{833}{4}
  \repere{1290}{856}{5}}

\begin{tabularx}{\linewidth}{@{}c X@{}}
\rep{1} & \ui{Individu} pour un animal seul, \ui{Troupeau / groupe d'animaux} pour un groupe. \\
\rep{2} & L'espèce : elle filtre les référentiels et les courbes de croissance proposés. \\
\rep{3} & La date de naissance : saisie directe, \ui{Sélectionner une date} ou \ui{Aujourd'hui}. \\
\rep{4} & Le sexe et le statut reproducteur. \\
\rep{5} & Le bouton coché enregistre le dossier. \\
\end{tabularx}

\begin{pourquoi}
L'espèce est déterminante : elle filtre les référentiels et les courbes de croissance
proposés. Le sexe et la date de naissance servent aux courbes de croissance et au calcul de
l'âge. Pour un groupe d'animaux, choisissez plutôt \ui{Troupeau} (section~\ref{sec:troupeau}).
\end{pourquoi}

% ------------------------------------------------------------------
\etape{Étape 2 — Ouvrir la fiche de l'animal}{repérer les onglets de travail.}

Dans la liste, cliquez sur la flèche de la ligne de Filou. La fiche s'ouvre sur l'onglet
\ui{Identification}.

\captureannotee{08_fiche_identification.png}{La fiche de l'animal.}{fiche}{%
  \repere{150}{40}{1}
  \repere{222}{116}{2}
  \repere{222}{172}{3}
  \repere{222}{228}{4}
  \repere{222}{312}{5}
  \repere{222}{396}{6}
  \repere{222}{649}{7}
  \repere{222}{781}{8}
  \repere{1336}{72}{9}}

\begin{tabularx}{\linewidth}{@{}c X@{}}
\rep{1} & Nom, espèce et identifiant de l'animal ouvert. \\
\rep{2} & \ui{Identification} : les informations saisies à l'étape 1. \\
\rep{3} & \ui{Consultations} : l'historique des consultations (étape 3). \\
\rep{4} & \ui{Rations} : composition et analyse des rations (étapes 5 à 7). \\
\rep{5} & \ui{Graphique} et \ui{Graphique Aliments} : suivi du poids et comparaisons (section~\ref{sec:graphiques}). \\
\rep{6} & \ui{Ordonnance} : documents pour le propriétaire (étape 8). \\
\rep{7} & \ui{Exporter animal} et \ui{Partager en ligne} : transmettre le dossier. \\
\rep{8} & \ui{Retour} à la liste des animaux. \\
\rep{9} & Modifier ou supprimer le dossier. \\
\end{tabularx}

% ------------------------------------------------------------------
\etape{Étape 3 — Saisir la consultation}{décrire l'état de Filou aujourd'hui.}

\begin{enumerate}
  \item Ouvrez l'onglet \ui{Consultations}, puis cliquez sur \ui{+} en bas à droite.
  \item Vérifiez la \ui{Date} \rep{1} et saisissez le \ui{Poids} \rep{2} : 32.
  \item Cliquez sur les trois points de la \ui{Note d'état corporel} \rep{3} et choisissez
        \emph{7/9 — Surpoids} (figure~\ref{fig:nec}).
  \item Lisez l'\ui{Estimation} du poids idéal \rep{5} : 26,67~kg. Si vous la retenez,
        saisissez-la dans \ui{Poids idéal (kg)} \rep{4}.
  \item Saisissez l'\ui{Objectif} \rep{6} : \emph{Perte de poids progressive}. Ajoutez si besoin
        des \ui{Mots-clés} \rep{7} (pathologies, contexte) et une observation.
\end{enumerate}
Ne validez pas encore : il reste à choisir le référentiel (étape 4).

\captureannotee{10_consultation_edition.png}{Saisie de la consultation.}{consultation}{%
  \repere{262}{100}{1}
  \repere{262}{180}{2}
  \repere{1328}{335}{3}
  \repere{262}{464}{4}
  \repere{1104}{447}{5}
  \repere{262}{551}{6}
  \repere{430}{649}{7}
  \repere{1356}{805}{8}}

\rep{8} Le bouton coché enregistre la consultation (étape 4).

\capture{11_consultation_nec.png}{Choix de la note d'état corporel sur 9.}{nec}

\begin{pourquoi}
Le poids idéal est le poids sur lequel les besoins sont calculés. Pour un animal en surpoids,
calculer les besoins sur le poids actuel reviendrait à entretenir l'excès de poids.
\begin{itemize}
  \item L'\ui{Estimation} est une suggestion calculée d'après le poids et la NEC
        (section~\ref{sec:calculs}).
  \item Elle n'est utilisée que si vous la recopiez dans \ui{Poids idéal (kg)}. Sinon, les
        calculs portent sur le poids actuel.
\end{itemize}
\end{pourquoi}

% ------------------------------------------------------------------
\etape{Étape 4 — Choisir le référentiel}{dire à l'application à quels besoins comparer la ration.}

\begin{enumerate}
  \item Faites défiler le formulaire jusqu'à \ui{Références nutritionnelles}.
  \item Dans \ui{Référence générale}, cliquez sur les trois points : la liste des référentiels
        de l'espèce s'ouvre.
  \item Tapez quelques lettres dans la recherche \rep{1} si la liste est longue, puis cliquez
        sur la ligne voulue \rep{3} : ici \emph{Adulte >25kg}.
  \item Si l'animal a une maladie à prendre en compte, ajoutez-la dans
        \ui{Références complémentaires}.
  \item Validez la consultation avec le bouton coché.
\end{enumerate}

\captureannotee{13_choix_referentiel.png}{Choix du référentiel général.}{referentiel}{%
  \repere{418}{263}{1}
  \repere{418}{306}{2}
  \repere{418}{462}{3}}

\rep{2} Le nombre de référentiels qui correspondent à la recherche.

\begin{pourquoi}
Le référentiel contient les besoins (minimum, optimum, maximum) de chaque nutriment pour un type
d'animal, ainsi que l'équation du besoin énergétique. C'est la règle à laquelle la ration sera
comparée : un mauvais référentiel donne une analyse fausse, même avec une bonne ration.
Les \ui{Coefficients d'ajustement énergétique} K1 à K5, plus bas dans le formulaire, modulent
le besoin énergétique (activité, stérilisation, environnement\ldots).
\end{pourquoi}

\begin{attention}
Une consultation sans référence générale ne peut pas être enregistrée : l'application affiche
\ui{Informations manquantes}.
\end{attention}

\begin{verifiez}
La consultation apparaît dans la liste avec la date, le poids et \ui{Rations: 1}.
L'application a créé une première ration vide, la \og Ration actuelle\fg.
\end{verifiez}

% ------------------------------------------------------------------
\etape{Étape 5 — Composer la ration}{mettre un aliment dans la ration et fixer sa quantité.}

\begin{enumerate}
  \item Ouvrez l'onglet \ui{Rations}. Dans le cadre \ui{Aliments de la ration}, cliquez sur
        \ui{+}.
  \item Tapez un nom, une marque ou un ingrédient \rep{1} : ici \emph{satiety}. Affinez si
        besoin avec les filtres \rep{2} (type, indications, espèce, base de données).
  \item Cliquez sur un aliment de la liste \rep{3} : sa fiche s'affiche à droite \rep{4}
        (marque, espèces, indications, ingrédients, composition). Vérifiez qu'il s'agit bien
        d'un aliment pour chien.
  \item Cliquez sur le bouton \ui{+} rond \rep{5} : l'aliment est ajouté avec une quantité de
        100~g. Le bouton \ui{+R} \rep{6} l'ajoute plutôt dans une nouvelle ration.
  \item Revenez à la ration avec la flèche en haut à gauche.
\end{enumerate}

\captureannotee{14_ajout_aliment.png}{Recherche et ajout d'un aliment.}{ajout-aliment}{%
  \repere{270}{122}{1}
  \repere{270}{212}{2}
  \repere{258}{605}{3}
  \repere{940}{124}{4}
  \repere{1295}{768}{5}
  \repere{1295}{848}{6}}

Pour régler la quantité, cliquez sur le crayon de la ligne de l'aliment, saisissez la quantité
en grammes (ici 380) et validez par \ui{OK}.

% ------------------------------------------------------------------
\etape{Étape 6 — Lire l'analyse}{savoir si la ration couvre les besoins de Filou.}

Avec 380~g de l'aliment, l'onglet \ui{Rations} se présente ainsi.

\captureannotee{15_ration_avant_ajustement.png}{Analyse de la ration de 380~g.}{ration-avant}{%
  \repere{252}{103}{1}
  \repere{636}{103}{2}
  \repere{1020}{101}{3}
  \repere{252}{212}{4}
  \repere{252}{576}{5}
  \repere{836}{212}{6}
  \repere{1180}{185}{7}
  \repere{836}{248}{8}}

\begin{tabularx}{\linewidth}{@{}c X@{}}
\rep{1} & \ui{Valeurs métaboliques} : poids, poids idéal, poids métabolique, besoin énergétique
          standard (BEE) et besoin énergétique (BE). \\
\rep{2} & \ui{Coefficients} K1 à K5 et coefficient d'ajustement. Le crayon les modifie. \\
\rep{3} & \ui{Bilan énergétique} : énergie apportée, couverture du besoin, K observé et K
          calculé. \textbf{C'est le premier chiffre à regarder.} \\
\rep{4} & Les rations de la consultation. \\
\rep{5} & Les aliments de la ration sélectionnée et les outils d'ajustement. \\
\rep{6} & \ui{Apports} : chaque nutriment, exprimé dans l'unité du référentiel (ici \og par
          1000 kcal\fg). \\
\rep{7} & Unité d'expression, nutriments affichés, export PDF, vue en cartes ou en graphique. \\
\rep{8} & Les nutriments par famille. Une icône d'alerte signale un écart au référentiel. \\
\end{tabularx}

\subsubsection{Interpréter le résultat}

La couverture énergétique n'est que de 56~\% : 380~g apportent 971~kcal/j pour un besoin de
1749~kcal/j. Presque tous les nutriments sont signalés en carence.

\begin{pourquoi}
Les apports sont rapportés au besoin énergétique de l'animal, et non à l'énergie de la ration
(section~\ref{sec:calculs}). Tant que la ration ne couvre pas le besoin énergétique, tous les
apports paraissent bas. \textbf{Corrigez donc toujours l'énergie en premier}, puis regardez les
autres nutriments.
\end{pourquoi}

% ------------------------------------------------------------------
\etape{Étape 7 — Ajuster la ration}{obtenir une couverture énergétique proche de 100~\%.}

\subsubsection{L'ajustement rapide}

Dans le cadre \ui{Aliments de la ration}, cliquez sur l'icône \ui{Ajustement rapide
multi-nutriments} (la balance, à gauche du livre et du \ui{+}). Les quantités sont recalculées
pour couvrir le besoin énergétique.

\captureannotee{17_ration_apres_ajustement.png}{La même ration après ajustement rapide.}{ration-apres}{%
  \repere{1160}{150}{1}
  \repere{252}{646}{2}
  \repere{1000}{406}{3}}

\begin{verifiez}
\rep{1} La couverture passe à 99~\% et le K observé à 0,99. \rep{2} La quantité passe à
675~g. \rep{3} Les alertes de carence disparaissent ; de petites flèches rouges signalent
désormais quelques apports au-dessus de l'optimum maximum (calcium, phosphore, vitamine~A).
\end{verifiez}

\begin{pourquoi}
Un dépassement de l'optimum n'est pas forcément un problème : survolez l'icône pour lire le
seuil concerné et jugez-en cliniquement. Pour une perte de poids, l'énergie se règle par le
poids idéal et les coefficients K, pas en donnant moins que la ration calculée : sinon, les
autres nutriments deviennent insuffisants.
\end{pourquoi}

\subsubsection{L'ajustement multi-nutriments}

L'icône voisine ouvre l'outil complet, utile quand la ration associe plusieurs aliments.

\captureannotee{16_ajustement_multinutriments.png}{L'ajustement multi-nutriments.}{ajustement}{%
  \repere{268}{147}{1}
  \repere{268}{320}{2}
  \repere{1385}{240}{3}
  \repere{1220}{448}{4}
  \repere{268}{552}{5}
  \repere{268}{719}{6}}

\begin{tabularx}{\linewidth}{@{}c X@{}}
\rep{1} & Actions rapides : verrouiller ou déverrouiller toutes les quantités, ajustement
          automatique, réinitialisation. \\
\rep{2} & Pour chaque aliment, le \ui{Nutriment cible} qu'il doit couvrir. \\
\rep{3} & Le cadenas fige la quantité d'un aliment. \\
\rep{4} & Le niveau visé pour chaque nutriment cible (minimum, optimum\ldots). \\
\rep{5} & La part du besoin énergétique couverte par les aliments complets, quand la ration
          associe aliment complet et compléments. \\
\rep{6} & \ui{Prévisualiser} sans modifier, \ui{Ajuster}, \ui{Quantité égale (Sanalio)}, ou
          \ui{Ajustement par contraintes (min/max)}, qui cherche à respecter toutes les bornes
          du référentiel à la fois. \\
\end{tabularx}

\subsubsection{Voir la ration en graphique}

L'interrupteur \ui{Cartes} / \ui{Bullet} du panneau \ui{Apports} affiche la composition de la
ration, l'origine de l'énergie et, pour chaque nutriment, une barre comparée aux seuils du
référentiel. \ui{Cliquer pour plus de détails} ouvre le détail d'un nutriment.

\captureannotee{18_analyse_graphique.png}{Analyse de la ration en vue graphique.}{analyse-graphique}{%
  \repere{1380}{240}{1}
  \repere{862}{256}{2}
  \repere{1135}{256}{3}
  \repere{1030}{700}{4}
  \repere{1190}{595}{5}}

\begin{tabularx}{\linewidth}{@{}c X@{}}
\rep{1} & L'interrupteur \ui{Cartes} / \ui{Bullet} bascule entre la vue en cartes et cette vue graphique. \\
\rep{2} & \ui{Composition} : la part de chaque constituant dans la ration, en masse. \\
\rep{3} & \ui{Origine Énergie} : la part de l'énergie apportée par les protéines, les lipides et l'ENA. \\
\rep{4} & Pour chaque nutriment, une barre situe l'apport par rapport aux seuils du référentiel. \\
\rep{5} & \ui{Cliquer pour plus de détails} ouvre le détail du nutriment. \\
\end{tabularx}

% ------------------------------------------------------------------
\etape{Étape 8 — Préparer l'ordonnance}{remettre au propriétaire une ration claire et des conseils.}

\begin{enumerate}
  \item Ouvrez l'onglet \ui{Ordonnance}.
  \item Cochez la ou les rations à prescrire \rep{1}.
  \item Ajoutez des conseils de la bibliothèque ou rédigez une \ui{Nouvelle section} \rep{2}.
  \item Complétez le compte rendu \rep{3} : anamnèse, examen clinique, facteur nutritionnel
        clé, texte additionnel.
  \item Utilisez les boutons du bas \rep{4} : \ui{Copier le compte rendu} (pour le coller dans
        votre logiciel), \ui{Aperçu du compte rendu}, \ui{Prévisualiser l'ordonnance}.
\end{enumerate}

\captureannotee{21_ordonnance.png}{Préparation de l'ordonnance et du compte rendu.}{ordonnance}{%
  \repere{278}{160}{1}
  \repere{268}{325}{2}
  \repere{830}{175}{3}
  \repere{585}{861}{4}}

\begin{aretenir}
Les coordonnées du praticien qui figurent sur l'ordonnance se saisissent une fois pour toutes
dans \ui{Paramètres} > \ui{Préférences} (section~\ref{sec:parametres}).
\end{aretenir}

% ------------------------------------------------------------------
\subsection{Récapitulatif du cas}

\begin{tabularx}{\linewidth}{@{}l l X@{}}
\toprule
\textbf{Étape} & \textbf{Où} & \textbf{Résultat pour Filou} \\
\midrule
1. Dossier & Liste, \ui{+} & Labrador mâle, né en 2019 \\
3. Consultation & \ui{Consultations} & 32 kg, NEC 7/9, poids idéal estimé 26,67 kg \\
4. Référentiel & formulaire de consultation & \emph{Adulte >25kg} \\
5. Ration & \ui{Rations}, \ui{+} & un aliment hypocalorique, 380 g \\
6. Analyse & \ui{Rations} & couverture 56 \%, carences apparentes \\
7. Ajustement & ajustement rapide & 675 g, couverture 99 \% \\
8. Ordonnance & \ui{Ordonnance} & ration cochée, compte rendu, impression \\
\bottomrule
\end{tabularx}

\begin{aretenir}
Lors d'un suivi, ne modifiez pas l'ancienne consultation : créez-en une nouvelle, ou dupliquez
la précédente avec l'icône de copie. L'historique des poids et des rations reste ainsi lisible
dans les graphiques.
\end{aretenir}

% ================================================================== 5
\section{Comprendre les calculs}
\label{sec:calculs}

Cette section explique les chiffres affichés dans l'onglet \ui{Rations}, avec les valeurs de
Filou.

\subsection{Le poids de calcul}

Les besoins sont calculés sur le \textbf{poids idéal} s'il est saisi, sinon sur le poids
actuel. L'estimation proposée dans la consultation considère que chaque point de NEC au-dessus
de 5 correspond à 10~\% de poids en trop :
\[
  \text{poids idéal estimé} = \frac{\text{poids} \times 100}{100 + (\text{NEC} - 5) \times 10}
  \qquad\text{Filou : } \frac{32 \times 100}{120} = 26{,}67\ \text{kg}
\]

\subsection{Du poids au besoin énergétique}

\begin{enumerate}
  \item Le \textbf{poids métabolique} tient compte du fait que les besoins n'augmentent pas
        proportionnellement au poids :
        \[ P_{\text{mét}} = P^{0{,}75} \qquad\text{Filou : } 32^{0{,}75} = 13{,}45\ \text{kg}^{0{,}75} \]
  \item Le \textbf{besoin énergétique standard} (BEE) applique l'équation du référentiel. Pour
        le référentiel de l'exemple :
        \[ \text{BEE} = 130 \times P^{0{,}75} = 130 \times 13{,}45 = 1749\ \text{kcal/j} \]
  \item Le \textbf{K calculé} est le produit des coefficients K1 à K5 et du coefficient
        d'ajustement. Le \textbf{besoin énergétique} (BE) en découle, augmenté le cas échéant de
        l'énergie demandée par les références de maladie :
        \[ \text{BE} = \text{BEE} \times K \;(+\ \text{énergie des références complémentaires}) \]
\end{enumerate}

\begin{aretenir}
Dans l'exemple, le poids idéal n'a pas été recopié : les calculs portent sur 32~kg. Avec
26,67~kg, le BEE tomberait à $130 \times 26{,}67^{0{,}75} \approx 1525$~kcal/j.
\end{aretenir}

\subsection{L'énergie apportée par les aliments}

L'énergie d'un aliment vient, dans l'ordre :
\begin{enumerate}
  \item de la valeur saisie pour l'espèce dans la fiche de l'aliment ;
  \item sinon de l'équation de densité énergétique du référentiel : pour le chien,
        l'équation des aliments commerciaux (aliments complets ou complémentaires), sinon
        l'équation d'Atwater (aliments bruts).
\end{enumerate}
L'équation des aliments commerciaux du chien tient compte des fibres :
\[
  \text{EM (kcal/100 g)} = \bigl(5{,}7\,P + 9{,}4\,L + 4{,}1\,(\text{ENA} + C)\bigr)
  \times \frac{91{,}2 - 1{,}43\,\dfrac{100\,C}{100 - H}}{100} - 1{,}04\,P
\]
avec $P$ les protéines, $L$ les lipides, $C$ la cellulose brute et $H$ l'humidité, en \%.
\begin{itemize}
  \item Si la cellulose brute n'est pas renseignée, la \textbf{fibre totale} la remplace.
  \item Si l'ENA n'est pas renseigné, il est calculé : $100 - P - H - C - \text{cendres} - L$.
        La cellulose approchée y est retirée, pour ne pas compter les fibres deux fois.
\end{itemize}
Exemple : un lait pour chiot à 33~\% de protéines, 39~\% de lipides, 19,5~\% d'ENA, 3~\%
d'humidité et 0,5~\% de fibre totale apporte 541,7~kcal/100~g, soit 1083~kcal pour 200~g.

\subsection{Le bilan énergétique}

\[
  \text{Couverture} = \frac{\text{énergie apportée}}{\text{BE}}
  \qquad\qquad
  K_{\text{observé}} = \frac{\text{énergie apportée}}{\text{BEE}}
\]
Avec 675~g : $1725 / 1749 = 99~\%$ et $K_{\text{observé}} = 0{,}99$. Le K observé indique quel
coefficient la ration applique réellement ; comparez-le au K calculé.

\subsection{Les apports}

Les apports sont exprimés dans l'unité du référentiel. En \og par 1000 kcal\fg, l'apport
journalier est rapporté au \textbf{besoin énergétique standard} (BEE), et non à l'énergie de la
ration :
\[
  \text{apport (g/1000 kcal)} = \frac{\text{apport journalier (g/j)}}{\text{BEE (kcal/j)}} \times 1000
\]
Les seuils du référentiel étant exprimés de la même façon, la comparaison revient à confronter
l'apport journalier au besoin journalier de l'animal.
Filou, avec 675~g d'un aliment à 30~\% de protéines : $202{,}5 / 1749 \times 1000 = 116$~g/1000~kcal.

\begin{aretenir}
Si la ration ne couvre que la moitié du besoin énergétique, chaque apport est divisé par deux
par rapport à ce qu'il serait avec la bonne quantité. D'où la règle : \textbf{l'énergie
d'abord}.
\end{aretenir}

\subsection{Les alertes}

\begin{tabularx}{\linewidth}{@{}>{\raggedright\arraybackslash}p{6.2cm} X@{}}
\toprule
\textbf{Message (au survol de l'icône)} & \textbf{Signification} \\
\midrule
Carence : apport inférieur au minimum & sous le seuil minimal du référentiel. \\
Carence : apport inférieur au optimal minimum & au-dessus du minimum, mais sous l'optimum. \\
Excès : apport supérieur au optimal maximum & au-dessus de l'optimum (flèche rouge). \\
Informations incomplètes & la composition d'au moins un aliment ne renseigne pas ce
nutriment : l'apport réel peut être plus élevé. \\
\bottomrule
\end{tabularx}

La section \ui{Analyses et ratios} donne les rapports K/Na, Ca/P, oméga-6/oméga-3 et Zn/Cu,
calculés sur la ration entière et arrondis au centième (Ca/P = 1,1 / 0,8 = 1,38).

% ================================================================== 6
\section{Aller plus loin avec les rations}

\subsection{Plusieurs rations dans une consultation}

Une consultation peut réunir la ration actuelle du propriétaire et une ou plusieurs
propositions, pour les comparer.
\begin{itemize}
  \item \ui{+} du cadre \ui{Rations de la consultation} ajoute une ration ; l'icône voisine
        reprend une ration d'une consultation précédente.
  \item Sur chaque ration : ajouter une étape au plan évolutif (transition alimentaire),
        dupliquer, modifier, supprimer.
  \item Dans \ui{Aliments de la ration} : enregistrer la ration comme recette, la copier, ouvrir
        les recettes enregistrées.
\end{itemize}

\subsection{Graphiques de suivi}
\label{sec:graphiques}

L'onglet \ui{Graphique} propose quatre vues : \ui{Évolution du poids}, \ui{Répartition
énergétique}, \ui{Densité énergétique des rations} et \ui{Nutriments des rations}.

\captureannotee{19_graphique_poids.png}{Évolution du poids et courbes de référence.}{graphique-poids}{%
  \repere{440}{176}{1}
  \repere{900}{176}{2}
  \repere{1095}{176}{3}
  \repere{895}{506}{4}
  \repere{265}{830}{5}
  \repere{745}{862}{6}
  \repere{1245}{857}{7}}

\begin{tabularx}{\linewidth}{@{}c X@{}}
\rep{1} & La courbe de référence proposée (sexe et classe de poids) ; un clic en choisit une autre. \\
\rep{2} & L'agrandissement des courbes de croissance et leur export en PDF. \\
\rep{3} & Afficher ou masquer les courbes de référence. \\
\rep{4} & Chaque pesée, placée selon l'âge. \\
\rep{5} & L'\ui{Historique des poids} de l'animal. \\
\rep{6} & L'objectif de perte de poids : perte minimale et maximale en \% par semaine, poids visé. \\
\rep{7} & \ui{+ Poids} ajoute une pesée. \\
\end{tabularx}


La vue du poids place chaque pesée en fonction de l'âge. Sous le graphique, l'\ui{Historique
des poids} permet d'ajouter une pesée et de fixer un objectif de perte de poids (perte minimale
et maximale en \% par semaine).

\begin{pourquoi}
La courbe de référence proposée par défaut correspond au sexe de l'animal et à sa classe de
poids. Cette classe est fixée, dans l'ordre :
\begin{enumerate}
  \item par le poids adulte attendu, pour un chiot (variable AW de la consultation) ;
  \item sinon par le poids idéal, saisi ou estimé d'après la NEC ;
  \item sinon par le dernier poids.
\end{enumerate}
Un animal en surpoids n'est ainsi pas classé trop haut. Le bouton au-dessus du graphique permet
de choisir une autre courbe.
\end{pourquoi}

Dans les vues de rations, chaque ration porte un numéro, repris dans la légende sous le
graphique : 1, 2\ldots{} pour les rations de la première consultation, 101, 102\ldots{} pour la
suivante. Deux options règlent l'affichage :
\begin{itemize}
  \item \ui{Consultation actuelle uniquement} limite le graphique à la consultation affichée ;
  \item \ui{Aire des possibilités et sommes} entoure les rations de chaque consultation
        (orange pour les rations actuelles, bleu pour les propositions) et ajoute leur moyenne
        pondérée par les coefficients, notée \textbf{$\Sigma$}. Les numéros des rations ne changent pas.
\end{itemize}

L'onglet \ui{Graphique Aliments} compare des aliments entre eux : ajoutez-les avec \ui{+} dans
la liste de gauche, puis cliquez sur \ui{Voir l'analyse graphique}.

\captureannotee{20_graphique_aliments.png}{Sélection d'aliments à comparer.}{graphique-aliments}{%
  \repere{268}{57}{1}
  \repere{495}{272}{2}
  \repere{752}{324}{3}
  \repere{1112}{272}{4}
  \repere{1125}{860}{5}}

\begin{tabularx}{\linewidth}{@{}c X@{}}
\rep{1} & La recherche et les filtres (type, espèce, base de données, indications). \\
\rep{2} & Les aliments trouvés ; le \ui{+} d'une ligne ajoute l'aliment à la comparaison. \\
\rep{3} & \ui{Tout ajouter} et \ui{Tout retirer}. \\
\rep{4} & Les aliments sélectionnés. \\
\rep{5} & \ui{Voir l'analyse graphique} ouvre les graphiques de comparaison. \\
\end{tabularx}

\subsection{Transmettre un dossier}

Depuis le menu de gauche de la fiche :
\begin{itemize}
  \item \ui{Exporter animal} enregistre le dossier complet (consultations, rations, aliments
        utilisés) dans un fichier, réimportable sur un autre poste ;
  \item \ui{Partager en ligne} transmet le dossier par un service en ligne.
\end{itemize}

% ================================================================== 7
\section{Cas particuliers}

\subsection{Un troupeau}
\label{sec:troupeau}

\textbf{Quand l'utiliser ?} Pour raisonner une ration à l'échelle d'un élevage ou d'un groupe :
un chenil, une chatterie, un lot d'animaux de rente.

À la création du dossier, choisissez \ui{Troupeau / groupe d'animaux}. Le troupeau n'a pas
d'espèce, de race ni de sexe propres : il est décrit par des groupes d'individus (nom,
effectif, poids moyen, espèce, race, sexe). Il peut mêler plusieurs espèces ; le référentiel et le
K de chaque groupe se choisissent dans chaque consultation.

\captureannotee{07_creation_troupeau.png}{Création d'un troupeau de douze individus.}{creation-troupeau}{%
  \repere{545}{152}{1}
  \repere{345}{441}{2}
  \repere{1105}{448}{3}
  \repere{474}{612}{4}
  \repere{305}{636}{5}
  \repere{1290}{856}{6}}

\begin{tabularx}{\linewidth}{@{}c X@{}}
\rep{1} & \ui{Troupeau / groupe d'animaux} est choisi. \\
\rep{2} & L'effectif total du troupeau, somme des groupes. \\
\rep{3} & Pour chaque groupe : son nom, le nombre d'individus et leur poids moyen ; la croix retire le groupe. \\
\rep{4} & L'espèce, la race et le sexe du groupe. \\
\rep{5} & \ui{Ajouter un groupe d'individus} ; un nouveau groupe reprend l'espèce du précédent. \\
\rep{6} & Le bouton coché enregistre le troupeau. \\
\end{tabularx}

\subsubsection{La consultation}

L'\ui{État du troupeau} remplace le poids, la NEC, le référentiel et les K de l'individu. Pour
chaque groupe, saisissez l'effectif et le poids moyen à la date de la consultation, la note
d'état corporel moyenne, le référentiel (de l'espèce du groupe) et le K du groupe, ainsi que les
variables éventuelles des équations du référentiel. Le besoin d'un individu vaut le besoin
standard du référentiel multiplié par le K du groupe.

\subsubsection{La ration et son analyse}

Les rations d'un troupeau sont saisies en quantités pour tout le troupeau, par jour. L'onglet
\ui{Rations} affiche au-dessus des rations le panneau du troupeau.

\captureannotee{37_troupeau_rations.png}{La ration d'un troupeau de douze chiens en deux groupes.}{troupeau-rations}{%
  \repere{1190}{128}{1}
  \repere{675}{145}{2}
  \repere{250}{168}{3}
  \repere{252}{275}{4}
  \repere{840}{398}{5}}

\begin{tabularx}{\linewidth}{@{}c X@{}}
\rep{1} & \ui{Ajuster la ration au besoin énergétique du troupeau} multiplie toutes les quantités
          par le facteur affiché (ici $\times 0{,}988$) pour que l'énergie apportée égale le besoin
          du troupeau. \\
\rep{2} & Ce qui est analysé : \ui{Troupeau entier}, ou un groupe (son nom et son effectif). \\
\rep{3} & La ration du troupeau par jour, le besoin du troupeau, l'énergie apportée et la
          couverture (verte de 90 à 110~\%, orange jusqu'à 80 et 120~\%, rouge au-delà). \\
\rep{4} & Une ligne par groupe : part de la ration, besoin et énergie d'un individu,
          couverture, nombre de nutriments hors normes et ration d'un individu. \\
\rep{5} & \ui{Apports au troupeau entier}, nutriment par nutriment. \\
\end{tabularx}

\begin{pourquoi}
La ration du troupeau est répartie entre les groupes au prorata de leur besoin énergétique :
un individu du groupe $i$ reçoit
\[
  \text{quantité} \times \frac{B_i}{\sum_j n_j\,B_j}
\]
où $n_j$ est l'effectif du groupe $j$ et $B_j$ le besoin énergétique d'un de ses individus. Dans
l'exemple, un adulte de 25~kg reçoit 271,6~g et un senior de 20~kg (K = 0,9) 206,8~g.
\end{pourquoi}

\textbf{Analyser un groupe.} Cliquez sur la ligne d'un groupe (ou sur sa puce) : la ration
affichée devient celle d'un individu de ce groupe, analysée comme celle d'un animal seul, avec
le poids, le référentiel et le K du groupe. Cette ration est en lecture seule ; cliquez à
nouveau sur la ligne, ou sur \ui{Troupeau entier}, pour revenir à la ration du troupeau et la
modifier.

\textbf{Le troupeau entier.} Le tableau \ui{Apports au troupeau entier} compare l'apport total
par jour à la somme des besoins des individus (effectif $\times$ besoin d'un individu, groupe par
groupe), pour chaque nutriment : énergie d'abord, puis les nutriments hors normes, puis les
autres par ordre alphabétique. \ui{Seulement les nutriments hors normes} réduit la liste.
Un astérisque signale une valeur incomplète : la composition d'au moins un aliment ne renseigne
pas ce nutriment.

\begin{attention}
Une somme conforme peut masquer la carence d'un groupe : un groupe trop nourri compense, dans le
total, un groupe qui manque. Regardez toujours la dernière colonne, \ui{Groupes hors normes}, et
le nombre d'écarts de chaque groupe dans le tableau du haut.
\end{attention}

\subsection{L'exploration multiration}
\label{sec:multiration}

\textbf{Quand l'utiliser ?} Pour préparer des rations types, par exemple pour toute une gamme
de poids, plutôt qu'une ration pour un animal précis.

Activez-la d'abord dans \ui{Paramètres} > \ui{Administration} > \ui{Activer l'analyse
multiration} ; le bouton \ui{Exploration multiration} apparaît dans la liste des animaux.

\captureannotee{28_multiration.png}{Configuration d'une exploration multiration.}{multiration}{%
  \repere{455}{150}{1}
  \repere{440}{191}{2}
  \repere{440}{250}{3}
  \repere{440}{343}{4}
  \repere{440}{724}{5}
  \repere{440}{815}{6}}

\begin{tabularx}{\linewidth}{@{}c X@{}}
\rep{1} & Le nombre de scénarios à calculer et ce qui empêche encore le calcul. \\
\rep{2} & \ui{Calculer toutes les rations}, actif une fois la configuration complète. \\
\rep{3} & Les onglets : \ui{Animaux}, \ui{Ingrédients}, \ui{Options}, \ui{Enregistrées}. \\
\rep{4} & L'espèce et les référentiels généraux à croiser. \\
\rep{5} & La grille des poids (de, à, pas). \\
\rep{6} & La grille du coefficient K (de, à, pas). \\
\end{tabularx}

\begin{enumerate}
  \item Choisissez l'espèce et cochez un ou plusieurs référentiels généraux.
  \item Définissez la grille : poids (de, à, pas) et coefficient K (de, à, pas). Le nombre de
        scénarios s'affiche.
  \item Dans l'onglet \ui{Ingrédients}, composez les listes d'aliments sources de protéines, de
        fibres, de calcium, d'oméga-6, de sel et d'énergie. L'ajustement suit cet ordre.
  \item Cliquez sur \ui{Calculer toutes les rations}. Un clic sur une case de la carte affiche
        ses rations, que l'on peut ouvrir dans l'analyse habituelle pour les modifier et les
        exporter.
\end{enumerate}
L'onglet \ui{Enregistrées} conserve les configurations nommées ; les résultats sont recalculés à
chaque ouverture.

% ================================================================== 8
\section{Gérer les données}

\subsection{La base d'aliments}
\label{sec:aliments}

\ui{Liste des aliments} (écran principal) affiche toute la base, avec les mêmes filtres que la
recherche d'aliments d'une ration. Le crayon modifie un aliment, la corbeille le supprime,
\ui{+} en crée un.

\captureannotee{22_liste_aliments.png}{Liste des aliments.}{aliments}{%
  \repere{690}{121}{1}
  \repere{700}{196}{2}
  \repere{72}{295}{3}
  \repere{780}{280}{4}
  \repere{1300}{856}{5}}

\begin{tabularx}{\linewidth}{@{}c X@{}}
\rep{1} & La recherche par nom, marque ou ingrédient. \\
\rep{2} & Les filtres : type, indications, espèce, base de données. \\
\rep{3} & Le crayon modifie l'aliment, la corbeille le supprime. \\
\rep{4} & Le type d'aliment, les espèces et les indications. \\
\rep{5} & \ui{+} crée un aliment. \\
\end{tabularx}

La fiche d'un aliment a trois onglets : \ui{Informations générales} (nom, marque, gamme, base
de données, prix, conditionnement, type), \ui{Composition nutritionnelle} et
\ui{Bibliographie}. Enregistrez avec \ui{Sauvegarder}.

\captureannotee{23_aliment_edition.png}{Modification d'un aliment.}{aliment-edition}{%
  \repere{615}{104}{1}
  \repere{480}{169}{2}
  \repere{1385}{500}{3}
  \repere{1385}{880}{4}}

\begin{tabularx}{\linewidth}{@{}c X@{}}
\rep{1} & \ui{Sauvegarder} enregistre les modifications. \\
\rep{2} & Les trois onglets : \ui{Informations générales}, \ui{Composition nutritionnelle}, \ui{Bibliographie}. \\
\rep{3} & La base de données d'origine de l'aliment. \\
\rep{4} & Le type d'aliment : complet ou complémentaire (énergie par l'équation des aliments commerciaux), sinon aliment brut (équation d'Atwater), voir section~\ref{sec:calculs}. \\
\end{tabularx}

\begin{aretenir}
Une composition incomplète produit des alertes \og Informations incomplètes\fg{} dans l'analyse.
Pour un aliment que vous créez, renseignez au moins l'énergie, les macronutriments, le calcium
et le phosphore.
\end{aretenir}

\subsection{Les données de calcul}
\label{sec:donnees}

Le bouton \ui{Données de calcul} donne accès aux éléments sur lesquels reposent les calculs.
Ils sont livrés avec l'application. Modifiez-les avec prudence : toute modification change les
analyses de toutes les consultations qui les utilisent.

\begin{tabularx}{\linewidth}{@{}l X@{}}
\toprule
\textbf{Onglet} & \textbf{Contenu} \\
\midrule
\ui{Équations} & formules du besoin énergétique, du poids métabolique et de la densité
énergétique (par exemple \texttt{130*BW\^{}0.75}) \\
\ui{Bibliographie} & références qui justifient les référentiels et les valeurs des aliments \\
\ui{Systèmes de calcul} & les référentiels : besoins par espèce, stade ou maladie \\
\ui{Conseils personnalisés} & bibliothèque de conseils reprise dans l'ordonnance \\
\bottomrule
\end{tabularx}

\captureannotee{26_referentiels.png}{Les référentiels nutritionnels (onglet Systèmes de calcul).}{referentiels}{%
  \repere{700}{56}{1}
  \repere{650}{156}{2}
  \repere{190}{257}{3}
  \repere{275}{489}{4}
  \repere{1222}{274}{5}
  \repere{1300}{856}{6}}

\begin{tabularx}{\linewidth}{@{}c X@{}}
\rep{1} & Les quatre onglets des données de calcul. \\
\rep{2} & La recherche d'un référentiel. \\
\rep{3} & Nom, espèce et stade physiologique du référentiel. \\
\rep{4} & Une ligne rouge signale un référentiel de maladie, à ajouter en référence complémentaire. \\
\rep{5} & Modifier, dupliquer ou supprimer le référentiel. \\
\rep{6} & \ui{+} crée un référentiel. \\
\end{tabularx}
\capture{24_equations.png}{Les équations.}{equations}
\capture{25_bibliographie.png}{La bibliographie.}{bibliographie}
\capture{27_conseils.png}{La bibliothèque de conseils.}{conseils}

% ================================================================== 9
\section{Paramètres}
\label{sec:parametres}

La roue dentée de l'écran principal ouvre les paramètres, en six onglets.

\begin{tabularx}{\linewidth}{@{}l X@{}}
\toprule
\textbf{Onglet} & \textbf{À quoi il sert} \\
\midrule
\ui{Interface} & taille de l'interface (50 à 200~\%) ; accès aux \ui{Analyses croisées} \\
\ui{Préférences} & coordonnées du praticien pour l'ordonnance ; unité d'expression des
besoins par espèce \\
\ui{Import/Export} & échanges de données (animaux, aliments, recettes, équations, références,
conseils) \\
\ui{Excel} & export et import des aliments en CSV compatible Excel \\
\ui{Recettes} & recettes enregistrées depuis les rations \\
\ui{Administration} & multiration, réimport des données initiales, sauvegardes, import
VetNutri~2, maintenance \\
\bottomrule
\end{tabularx}

\subsection{Préférences : à faire dès l'installation}

Saisissez vos informations (nom, numéro d'ordre, adresse, téléphone, e-mail), reprises sur
l'ordonnance, puis cliquez sur \ui{Enregistrer}. La section \ui{Expression des besoins par
espèce} fixe l'unité d'expression des besoins (par exemple \og par 1000 kcal\fg).

\captureannotee{30_param_preferences.png}{Préférences.}{param-preferences}{%
  \repere{700}{58}{1}
  \repere{290}{145}{2}
  \repere{190}{645}{3}
  \repere{395}{713}{4}
  \repere{1300}{812}{5}}

\begin{tabularx}{\linewidth}{@{}c X@{}}
\rep{1} & Les six onglets des paramètres. \\
\rep{2} & Les coordonnées du praticien, reprises sur l'ordonnance. \\
\rep{3} & \ui{Enregistrer} les coordonnées. \\
\rep{4} & \ui{Expression des besoins par espèce} : l'unité dans laquelle les besoins sont affichés. \\
\rep{5} & Cliquez sur une espèce pour changer son unité (par exemple \og par 1000 kcal\fg). \\
\end{tabularx}

\subsection{Interface}

\capture{29_param_interface.png}{Taille de l'interface.}{param-interface}

\subsection{Import, export et Excel}

L'onglet \ui{Import/Export} exporte et importe des données au format API, avec une sélection
avancée, et importe des aliments, des animaux ou des références nutritionnelles
(\texttt{.vbnr.json}).

Dans un fichier de références, chaque besoin indique son unité : par Mcal (par 1000 kcal), par
kg de poids, par kg de poids métabolique, en valeur absolue ou en ratio. Un besoin dont l'unité
n'est pas reconnue (par exemple \og DM\fg) n'est pas importé : il apparaît dans le compte rendu
de l'import, pour que vous corrigiez le fichier.

\capture{31_param_import_export.png}{Import et export de données.}{param-import}

L'onglet \ui{Excel} exporte et importe les aliments en CSV : séparateur point-virgule,
encodage UTF-8, une ligne d'en-têtes, une colonne par nutriment avec l'unité dans l'en-tête.
Une prévisualisation permet de vérifier les données avant de confirmer l'import.

\capture{32_param_excel.png}{Import et export Excel des aliments.}{param-excel}

\subsection{Recettes}

\capture{33_param_recettes.png}{Recettes enregistrées.}{param-recettes}

\subsection{Administration}

\captureannotee{34_param_administration.png}{Administration de la base.}{param-admin}{%
  \repere{1285}{272}{1}
  \repere{40}{338}{2}
  \repere{40}{402}{3}
  \repere{40}{466}{4}
  \repere{40}{530}{5}
  \repere{40}{603}{6}}

\begin{tabularx}{\linewidth}{@{}c X@{}}
\rep{1} & \ui{Activer l'analyse multiration} affiche le bouton \ui{Exploration multiration} dans la liste des animaux. \\
\rep{2} & Relancer l'import automatique des données initiales (aliments, référentiels, conseils). \\
\rep{3} & \ui{Gérer les sauvegardes} (section~\ref{sec:sauvegardes}). \\
\rep{4} & Importer les données d'une installation VetNutri~2. \\
\rep{5} & Ajouter une même référence bibliographique à de nombreux aliments. \\
\rep{6} & Les actions irréversibles : vider une base de données. \\
\end{tabularx}

\begin{attention}
Les boutons rouges \ui{Vider la base de données\ldots} suppriment définitivement les aliments,
les animaux, les référentiels, les équations ou les bibliographies. Faites une sauvegarde avant
de les utiliser.
\end{attention}

\subsection{Sauvegardes}
\label{sec:sauvegardes}

\ui{Administration} > \ui{Gérer les sauvegardes} liste les sauvegardes avec leur date, leur
contenu et leur taille. \ui{Nouvelle sauvegarde} en crée une ; sur chaque ligne, une icône la
restaure et une autre la supprime.

\captureannotee{35_sauvegardes.png}{Gestion des sauvegardes.}{sauvegardes}{%
  \repere{1140}{40}{1}
  \repere{320}{121}{2}
  \repere{700}{176}{3}
  \repere{1270}{120}{4}
  \repere{1375}{120}{5}}

\begin{tabularx}{\linewidth}{@{}c X@{}}
\rep{1} & \ui{Nouvelle sauvegarde}. \\
\rep{2} & La date et le fichier de la sauvegarde. \\
\rep{3} & Son contenu : animaux, aliments, équations, conseils, et sa taille. \\
\rep{4} & Restaurer cette sauvegarde : ses données sont fusionnées avec les données actuelles, après une sauvegarde automatique préalable. \\
\rep{5} & Supprimer la sauvegarde. \\
\end{tabularx}

\begin{aretenir}
Faites une sauvegarde avant toute opération importante : import massif, vidage d'une base,
changement d'ordinateur. Copiez de temps en temps une sauvegarde hors de l'ordinateur.
\end{aretenir}

\subsection{Analyses croisées}
\label{sec:analyses-croisees}

\textbf{Quand l'utiliser ?} Pour comparer de nombreuses consultations : travaux dirigés, séries
de cas, examens.

Filtrez par espèce, mots-clés, identifiant d'examen ou d'exercice, sélectionnez les
consultations, puis ouvrez l'analyse ou la notation.

\capture{36_analyses_croisees.png}{Analyses croisées.}{analyses-croisees}

% ================================================================== Annexes
\appendix

\section{Questions fréquentes}

\begin{description}[style=nextline,leftmargin=0pt,itemsep=6pt]
  \item[Je ne peux pas enregistrer la consultation.]
    Il manque la \ui{Référence générale}. Choisissez un référentiel dans la section
    \ui{Références nutritionnelles} du formulaire.
  \item[Aucun référentiel ne s'affiche dans la liste.]
    Vérifiez l'espèce de l'animal : seuls les référentiels de son espèce sont proposés. Vérifiez
    aussi que les données ont été installées (écran d'accueil).
  \item[Tous les nutriments sont en carence.]
    Regardez la couverture énergétique : si elle est loin de 100~\%, ajustez d'abord la
    quantité (section~\ref{sec:calculs}).
  \item[Le besoin énergétique me paraît trop élevé pour un animal en surpoids.]
    Les calculs utilisent le poids actuel tant que le poids idéal n'est pas saisi. Recopiez
    l'estimation, ou votre propre valeur, dans \ui{Poids idéal (kg)}.
  \item[La courbe de croissance ne correspond pas à mon chiot.]
    Saisissez son poids adulte attendu (variable AW) dans la consultation, ou choisissez une
    autre courbe avec le bouton au-dessus du graphique.
  \item[Le bouton \ui{Exploration multiration} n'apparaît pas.]
    Activez-le dans \ui{Paramètres} > \ui{Administration}.
  \item[L'énergie d'un aliment riche en fibres me paraît basse.]
    Sans cellulose brute renseignée, la fibre totale, souvent plus élevée, est utilisée à sa
    place dans l'équation (section~\ref{sec:calculs}). Saisissez la cellulose brute dans la
    fiche de l'aliment si vous la connaissez.
  \item[Un aliment manque dans la base.]
    Créez-le depuis la \ui{Liste des aliments}, ou directement depuis la recherche d'aliments
    d'une ration (\ui{+} en haut à droite).
\end{description}

\section{Glossaire}

\begin{tabularx}{\linewidth}{@{}l X@{}}
\toprule
\textbf{Terme} & \textbf{Définition} \\
\midrule
BEE & Besoin énergétique standard, calculé par l'équation du référentiel (kcal/j). \\
BE & Besoin énergétique de l'animal : BEE × K, plus l'énergie liée aux maladies. \\
Cellulose brute & Fraction des fibres utilisée par l'équation d'énergie ; à défaut, la fibre totale. \\
Couverture & Part du BE apportée par la ration (\%). \\
ENA & Extractif non azoté (glucides digestibles), en \% ; calculé par différence s'il manque. \\
K, K1 à K5 & Coefficients qui modulent le besoin énergétique (activité, stérilisation\ldots). \\
K observé & Énergie apportée divisée par le BEE : le coefficient réellement appliqué. \\
NEC & Note d'état corporel, de 1 (très maigre) à 9 (très obèse) ; 4 et 5 correspondent à l'état idéal. \\
Poids idéal & Poids sur lequel les besoins sont calculés ; à défaut, le poids actuel. \\
Poids métabolique & Poids élevé à la puissance 0,75 (kg$^{0,75}$). \\
Référentiel & Ensemble des besoins (minimum, optimum, maximum) d'un type d'animal. \\
Ration actuelle & Première ration d'une consultation, souvent ce que mange déjà l'animal. \\
Troupeau & Dossier qui décrit un groupe d'animaux par groupes d'individus. \\
\bottomrule
\end{tabularx}

\section{Mémo : une consultation en dix gestes}

\begin{enumerate}
  \item Liste des animaux : \ui{+}, ou flèche d'un animal existant.
  \item \ui{Consultations} : \ui{+}, ou copie de la consultation précédente.
  \item Poids et NEC.
  \item Poids idéal : recopier l'estimation ou saisir sa valeur.
  \item Référence générale, et complémentaires en cas de maladie ; valider.
  \item \ui{Rations} : \ui{+} dans \ui{Aliments de la ration}, rechercher, ajouter.
  \item Lire le bilan énergétique.
  \item Ajustement rapide, ou ajustement multi-nutriments.
  \item Relire les alertes nutriment par nutriment.
  \item \ui{Ordonnance} : cocher la ration, ajouter conseils et compte rendu, imprimer.
\end{enumerate}

\end{document}
